%%%%%%%%%%  scan idx 199  |  folio 186  %%%%%%%%%%
\textbf{Lemme 2.4.8.2.} (cf. EGA IV 18.10.3) Soient $Y$ un espace pseudo-analytique normal\footnote{Il suffit que les anneaux locaux de $Y$ soient intègres.}, et $f:X\longrightarrow Y$ un morphisme non ramifié. L'ensemble des points de $X$ où $f$ est étale est à la fois ouvert et fermé.

\emph{Preuve.} Il est immédiat que l'ensemble des points où $f$ est étale est ouvert (car dire que $f$ est étale en un point signifie que $f$ est un isomorphisme au voisinage de ce point ([8] VI)). Il reste à montrer que l'ensemble des points où $f$ n'est pas étale est ouvert. Comme un morphisme non ramifié est localement une immersion ([8] VI), on peut supposer que $f$ est une immersion fermée, définie par un Idéal cohérent $I$. Soit $x$ un point de $X$ tel que $I_x\ne0$. Soit $\underline{p}$ un idéal premier de $\mathcal O_{Y,x}$ contenant $I_x$. On a $\dim(\mathcal O_{Y,x})_{\underline{p}}>0$, car dans le cas contraire on aurait $\underline{p}=0$ puisque $\mathcal O_{Y,x}$ est intègre. Donc $Y$ est de codimension $>0$ en $x$, donc de codimension $>0$ dans un voisinage de $x$ ([10] III A Cor 1 du lemme 3), et l'on gagne.

\section*{3. Compléments sur les faisceaux quasi-cohérents sur les schémas\protect\footnote{Ce numéro est dû essentiellement à J.L. Verdier.}}

\textbf{3.0.} Dans cette section, on se fixe un schéma \emph{quasi-compact et quasi-séparé} $S$. Le foncteur d'inclusion canonique (pleinement fidèle)
\[
 \phi:\Qcoh(S)\longrightarrow\Mod(S)\qquad\text{\footnotemark}
\]\footnotetext{$\Mod(S)$ (resp. $\Qcoh(S)$) désigne la catégorie des $\mathcal O_S$-Modules (resp. $\mathcal O_S$-Modules quasi-cohérents).}
s'étend en un foncteur, encore noté $\phi$, de $\mathrm D(\Qcoh(S))$ dans $\mathrm D(S)=\mathrm D(\Mod(S))$, qu'on se propose d'étudier\footnote{Nous omettrons de noter $\phi$ dans les formules où il n'y aura pas de confusion à craindre.}.

\textbf{Lemme 3.1.} La catégorie $\Qcoh(S)$ est une catégorie abélienne vérifiant (AB 5) et pour laquelle les faisceaux quasi-cohérents de type fini forment une «~petite~» famille de générateurs. En particulier (Tohoku 1.10.1), $\Qcoh(S)$ possède assez d'injectifs.

